%=============================================================================!
% 13.F  SOLUTIONS TO THE EXERCISES                                            !
%=============================================================================!
\unnumbered{Chapter~\ref{ch:thermostats}: Thermostats}
\label{sec:th-solutions}

\solutionto{ex:th-collisions}The intervals are independent, so the
probability of no collision in any of them is $(1 - \nu\dt/n)^n$, whose
logarithm $n\ln(1 - \nu\dt/n)$ tends to $-\nu\dt$ as $n$ grows, by $\ln(1 +
y) \approx y$ (\cref{sec:th-langevin}): the probability tends to
$\mathrm{e}^{-\nu\dt}$. At least one collision then has the probability
$1 - \mathrm{e}^{-0.002} = 0.001998$.

\solutionto{ex:th-lambda}$\lambda^2 = 1 + (1/500)(300/290 - 1) =
1.00006897$, so $\lambda = 1.00003448$: the velocities grow by 34 parts in
a million.

\solutionto{ex:th-mean}Each step adds the heat $(\dt/\tau_{\mathrm
T})(K_0 - K)$ (\cref{sec:th-berendsen}). Its average over the settled run
is zero, so $\ens{K} = K_0$, and since the kinetic temperature is
proportional to $K$, $\ens{T} = T_0$.

\solutionto{ex:th-langevin-tau}Averaging the square of \cref{eq:th-ou} with
$\sigma^2 = 2\gamma\kB T_0/m$ (\cref{eq:th-fdr}), the noise having mean 0
and variance 1 and being independent of $v$, $\ens{v'^2} = c^2\ens{v^2} + (1 -
c^2)\kB T_0/m$, so each component's kinetic energy changes over the step
by $(1 - c^2)\bigl(\frac12\kB T_0 - \frac12m\ens{v^2}\bigr)$. For a short
step $1 - c^2 = 1 - \mathrm{e}^{-2\gamma\dt} \approx 2\gamma\dt$, so the
rate is $2\gamma(\frac12\kB T_0 - \frac12m\ens{v^2})$, and over the $3N$
components, all of which the noise reaches, $2\gamma(K_0 - K)$ with $K_0 =
\frac32N\kB T_0$; the three freedoms of the centre of mass change $C_K$ by
\SI{0.4}{\percent}, which is negligible here. As in
\cref{sec:th-berendsen}, $C_V\,\dd T/\dd t = -2\gamma C_K(T - T_0)$, with
the time constant $C_V/(2\gamma C_K) = 1.590/(2\times\SI{1}{\per\pico
\second}) = \SI{0.795}{\pico\second}$; the fit in
\cref{fig:th-berendsen}(b) gives \SI{0.68}{\pico\second}, uncertain by
\SI{0.22}{\pico\second} in a single run, so the two agree.

\solutionto{ex:th-ice}$\ens{(\delta T/T_0)^2}/\tau_{\mathrm T} =
\SI{0.005}{\per\pico\second}$, so in \SI{1000}{\pico\second} the factor is
$\mathrm{e}^5 = 148$.

\solutionto{ex:th-gauss-sum}The density of $s$ adds up, over every $x$,
the chance that the first number is $x$ and the second $s - x$:
\begin{equation*}
   \mathcal{P}(s) = \frac{1}{2\pi\sqrt{ab}}\int\mathrm{e}^{-x^2/2a
   - (s - x)^2/2b}\,\dd x .
\end{equation*}
The exponent is $-[(a + b)x^2 - 2asx + as^2]/2ab$, and completing the
square,
\begin{equation*}
   (a + b)x^2 - 2asx + as^2 = (a + b)\Bigl(x - \frac{as}{a + b}\Bigr)^2
   + \frac{ab\,s^2}{a + b} ,
\end{equation*}
so the exponent is $-(a + b)(x - x_0)^2/2ab - s^2/2(a + b)$, with $x_0 =
as/(a + b)$. The Gaussian integral over $x$ gives $\sqrt{2\pi ab/(a +
b)}$, and
\begin{equation*}
   \mathcal{P}(s) = \frac{1}{\sqrt{2\pi(a + b)}}\,\mathrm{e}^{-s^2/2(a + b)} ,
\end{equation*}
the Gaussian of variance $a + b$.

\solutionto{ex:th-ou-variance}The two terms of \cref{eq:th-ou} are
independent, so their variances add, and multiplying a quantity by $c$
multiplies its variance by $c^2$: the variance after the step is
$c^2\sigma_v^2 + (\sigma^2/2\gamma)(1 - c^2)$, with $\sigma_v^2$ the
variance of $v$ before. It equals $\sigma_v^2$ only if $\sigma_v^2(1 -
c^2) = (\sigma^2/2\gamma)(1 - c^2)$, that is, since $c < 1$, if
$\sigma_v^2 = \sigma^2/2\gamma$.

\solutionto{ex:th-noise}With $c = \mathrm{e}^{-0.01} = 0.99005$,
\begin{equation*}
   \sqrt{1 - c^2}\,\sqrt{\kB T/m} = 0.14072\times
   \SI{0.001676}{\angstrom\per\femto\second}
   = \SI{2.36e-4}{\angstrom\per\femto\second} ,
\end{equation*}
a seventh of the thermal spread.

\solutionto{ex:th-baoab-vv}With $\gamma = 0$, $c = 1$ and the noise
vanishes, so O leaves the velocities unchanged; the two half drifts
then make one drift, and B A A B is a half kick, a drift and a half
kick, \cref{eq:in-kdk}.

\solutionto{ex:th-flux}With $w = mv^2/2\kB T$, $\dd w = mv\,\dd v/\kB T$,
so
\begin{equation*}
   \int_0^\infty v\,\mathrm{e}^{-mv^2/2\kB T}\dd v
   = \frac{\kB T}{m}\int_0^\infty\mathrm{e}^{-w}\dd w = \frac{\kB T}{m} .
\end{equation*}
Dividing by $\sqrt{2\pi\kB T/m}$,
with $\kB T/m = \sqrt{\kB T/m}\,\sqrt{\kB T/m}$, gives $\sqrt{\kB T/m}/
\sqrt{2\pi} = \sqrt{\kB T/2\pi m}$. For lithium at \SI{1000}{\kelvin},
with the factor 103.6427 of \cref{eq:en-mv2},
\begin{equation*}
   \sqrt{\frac{0.086173}{2\pi\times6.94\times103.6427}}
   = \SI{0.00437}{\angstrom\per\femto\second} .
\end{equation*}

\solutionto{ex:th-tst-1d}In one dimension the edge is a point and the
cell an interval, so $k_{\mathrm{TST}} = \sqrt{\kB T/2\pi m}\,
\mathrm{e}^{-U_{\mathrm b}/\kB T}/\int\mathrm{e}^{-U/\kB T}\dd x$. When
$U_{\mathrm b} \gg \kB T$ the integral comes almost wholly from the bottom
of the well, where it is the Gaussian integral $\int\mathrm{e}^{-m\omega^2
x^2/2\kB T}\dd x = \sqrt{2\pi\kB T/m\omega^2}$; dividing,
\begin{equation*}
   k_{\mathrm{TST}} = \sqrt{\frac{\kB T}{2\pi m}}\sqrt{\frac{m\omega^2}{2\pi
   \kB T}}\,\mathrm{e}^{-U_{\mathrm b}/\kB T}
   = \frac{\omega}{2\pi}\,\mathrm{e}^{-U_{\mathrm b}/\kB T} ,
\end{equation*}
the frequency of attempts times the Boltzmann factor of the barrier. With
$\omega/2\pi = 1/\SI{170.3}{\femto\second}$ and $\mathrm{e}^{-0.3/0.025852}
= \num{9.12e-6}$, $k_{\mathrm{TST}} = \SI{5.36e-8}{\per\femto\second}$, or
\num{5.4e7} hops per second. It is eleven times below the rate of
\cref{eq:th-tst} for the surface at \SI{300}{\kelvin}, \num{5.9e8} per
second: a hollow has six bridges to leave by, and across the path each
bridge has a third of the hollow's curvature, which widens the way out by
a further $\sqrt3$, and $6\sqrt3 = 10.4$.

\solutionto{ex:th-nh-swing}Without forces $\dot K = -2\xi K$, so
$\ddot K = -2\dot\xi K - 2\xi\dot K$. For small departures $\xi$ is small
and $K \approx K_0$, so to first order $\dd^2\,\delta K/\dd t^2 =
-2K_0\dot\xi = -2K_0\cdot2\,\delta K/Q$, since $\dot\xi = (2K -
2K_0)/Q$. With $K_0 = \frac12N_{\mathrm f}\kB T_0$ and $Q =
N_{\mathrm f}\kB T_0\tau_{\mathrm T}^2$, $4K_0/Q = 2/\tau_{\mathrm T}^2$. The
swing has the angular frequency $\sqrt2/\tau_{\mathrm T}$ and the period
$2\pi\tau_{\mathrm T}/\sqrt2 = \SI{444}{\femto\second}$ for $\tau_{\mathrm
T} = \SI{100}{\femto\second}$.

\solutionto{ex:th-nose}Hamilton's equations in $t'$ are $\dd\vb{r}_i/\dd
t' = \tilde{\vb{p}}_i/m_is^2$, $\dd\tilde{\vb{p}}_i/\dd t' = \vb{F}_i$,
$\dd s/\dd t' = p_s/Q$ and $\dd p_s/\dd t' = \sum_i\tilde{\vb{p}}_i^2/m_i
s^3 - N_{\mathrm f}\kB T_0/s$. With $\dd/\dd t = s\,\dd/\dd t'$:
\begin{align*}
   \frac{\dd\vb{r}_i}{\dd t} &= \frac{\tilde{\vb{p}}_i}{m_is}
   = \frac{\vb{p}_i}{m_i} ,\\
   \frac{\dd\vb{p}_i}{\dd t} &= s\,\frac{\dd}{\dd t'}\frac{\tilde{
   \vb{p}}_i}{s} = \frac{\dd\tilde{\vb{p}}_i}{\dd t'} - \frac{\tilde{
   \vb{p}}_i}{s}\frac{\dd s}{\dd t'} = \vb{F}_i - \xi\vb{p}_i ,\\
   \frac{\dd\xi}{\dd t} &= \frac{s}{Q}\frac{\dd p_s}{\dd t'}
   = \frac{s}{Q}\Bigl(\sum_i\frac{\tilde{\vb{p}}_i^2}{m_is^3}
   - \frac{N_{\mathrm f}\kB T_0}{s}\Bigr)
   = \frac{1}{Q}\Bigl(\sum_i\frac{\vb{p}_i^2}{m_i} - N_{\mathrm f}\kB
   T_0\Bigr) = \frac{2K - N_{\mathrm f}\kB T_0}{Q} ,
\end{align*}
the second line by the product rule, with $\xi = p_s/Q$.

\solutionto{ex:th-divergence}$\dot x = p_x/m$ does not depend on $x$, and
likewise for $y$ and $z$: three zeros. $\dot p_x = F_x - \xi p_x$ gives
$\partial\dot p_x/\partial p_x = -\xi$, and likewise for $p_y$ and $p_z$.
$\dot\xi = (2K - 3\kB T_0)/Q$ does not depend on $\xi$. The divergence is
$-3\xi$.

\solutionto{ex:th-csvr-mean}$R_1$ has mean 0 and $R_1^2$ mean 1, and $R_\perp^2$
has the mean $N_{\mathrm f} - 1$, so $\ens{\lambda^2} = c_1 + c_2
N_{\mathrm f} = c_1 + (1 - c_1)K_0/K$. Then $\ens{K'} = \ens{\lambda^2}K =
c_1K + (1 - c_1)K_0$, and $\ens{K'} - K = (1 - c_1)(K_0 - K)$.

\solutionto{ex:th-chi}In the canonical ensemble each mass-weighted
velocity component $u_k$ has a density proportional to
$\mathrm{e}^{-u_k^2/2\kB T}$ (\cref{sec:sm-maxwell}), which with $\kB T =
1$ is the normal distribution of variance 1; so $n$ independent normal
numbers are such components, and $K = \frac12\sum u_k^2$ has the density
\cref{eq:sm-gamma} with $N_{\mathrm f} = n$, proportional to $K^{n/2 -
1}\mathrm{e}^{-K}$: the gamma distribution of shape $n/2$. The sum of
squares is $2K$.

\solutionto{ex:th-trajcast}A hundred strides of \SI{5}{\femto\second} is
$\tau_{\mathrm T} = \SI{500}{\femto\second}$, and $c_1 = \mathrm{e}^{-5/500}
= 0.9900$; for paracetamol, ten strides of \SI{7}{\femto\second} make
\SI{70}{\femto\second}.

\solutionto{ex:th-com}$x_{\mathrm{com}} = \int_0^tv_{\mathrm{com}}(t_1)
\,\dd t_1$, and its square is the product of two copies,
$\int_0^tv_{\mathrm{com}}(t_1)\,\dd t_1\int_0^tv_{\mathrm{com}}(t_2)\,\dd
t_2$, a sum over every pair of times; the average of a sum is the sum of
the averages, so
\begin{align*}
   \ens{x_{\mathrm{com}}^2} &= \int_0^t\!\!\int_0^t
   \ens{v_{\mathrm{com}}(t_1)v_{\mathrm{com}}(t_2)}\,\dd t_1\dd t_2\\
   &= \frac{2\kB T}{M}\int_0^t\dd t_1\int_0^{t_1}\mathrm{e}^{-\gamma(t_1 -
   t_2)}\dd t_2
   && \text{pairs with $t_2 > t_1$ repeat those with $t_2 < t_1$}\\
   &= \frac{2\kB T}{M}\int_0^t\frac{1 - \mathrm{e}^{-\gamma t_1}}{\gamma}
   \,\dd t_1
   && \text{since $\textstyle\int_0^{t_1}\mathrm{e}^{\gamma t_2}\dd t_2 =
   (\mathrm{e}^{\gamma t_1} - 1)/\gamma$}\\
   &= \frac{2\kB T}{M\gamma}\Bigl(t - \frac{1 - \mathrm{e}^{-\gamma
   t}}{\gamma}\Bigr) .
\end{align*}
For $\gamma t \gg 1$ the bracket is $t - 1/\gamma \approx t$, so $D = \kB
T/(M\gamma)$; this is one component, and the three together give the $6Dt$
of \cref{sec:th-tests}. For 256 argon atoms, $M = \SI{10227}{\amu}$, with
$\kB T = \SI{0.011633}{\electronvolt}$ at \SI{135}{\kelvin}, $\gamma =
\SI{1e-4}{\per\femto\second}$ and the factor 103.6427 of \cref{eq:en-mv2}:
$0.011633/(10227\times103.6427\times10^{-4}) =
\SI{1.10e-4}{\angstrom\squared\per\femto\second}$.

\solutionto{ex:th-water}The true count is $3N - n_{\mathrm c} - 3 = 576 -
192 - 3 = 381$ (\cref{sec:sm-equipartition}). The thermostat drives
$2K/(576\kB)$ to \SI{300}{\kelvin}, while the true temperature is
$2K/(381\kB)$, larger by $576/381 = 1.512$: \SI{453.5}{\kelvin}.

\solutionto{ex:th-referee}Berendsen coupling with $\tau_{\mathrm T} =
\SI{0.1}{\pico\second}$ narrows the fluctuations of the kinetic energy
to 0.45 of the canonical spread, and those of the total energy with them,
while \cref{eq:sm-fluctuation} holds only for canonical fluctuations. The
heat capacity, which goes as the square of the energy's spread, comes out
too small: in this chapter's runs, $0.41N\kB$ against $2.376N\kB$. CSVR, a
Nosé-Hoover chain or Langevin dynamics give canonical fluctuations, and
with them 2.28 to $2.40N\kB$.

\solutionto{ex:th-diagnose}The noise does not keep the total momentum, so
the centre of mass diffuses, at $\kB T/(M\gamma) =
\SI{1.10e-4}{\angstrom\squared\per\femto\second}$ for this box over times
long against $1/\gamma = \SI{10}{\pico\second}$ (\cref{ex:th-com}), and
displacements measured in the cell include it. Over the fitted lags of 2
to \SI{20}{\pico\second}, $\gamma t$ is only 0.2 to 2, and the formula of
\cref{ex:th-com} rises there with 0.64 of its long-time slope, so the
centre of mass should add about $0.64\times1.10 = 0.70$: $3.69 + 0.70
\approx 4.4$, against the 4.62 measured, 1.2 times the 3.81 at fixed
energy. The centre of mass is one particle, so its share scatters widely,
0.55, 1.30 and 0.95 in the three runs. Measuring displacements from the
centre of mass gives 3.69, 0.97 of the value at fixed energy, as friction
requires.
